\documentclass[12pt,a4paper,UTF8]{ctexart}
\usepackage[a4paper,top=22mm,bottom=22mm,left=22mm,right=22mm,headheight=15pt]{geometry}
\newcommand{\Needspace}[1]{\par\begingroup\dimen0=#1\relax\dimen2=\pagegoal\advance\dimen2 by-\pagetotal\ifdim\dimen0>\dimen2\newpage\fi\endgroup}
\newcommand{\needspace}[1]{\par\begingroup\skip0=\lastskip\vskip-\skip0\dimen0=#1\relax\advance\dimen0 by\skip0\vskip\dimen0\penalty-100\vskip-\dimen0\vskip\skip0\penalty9999\endgroup}
\usepackage{amsmath,amssymb,mathtools,bm}
\usepackage{graphicx,xcolor,booktabs,longtable,array,enumitem,fancyhdr,caption,float,microtype,setspace}
\usepackage[hidelinks,unicode]{hyperref}
\usepackage[most]{tcolorbox}
\definecolor{ContestBlue}{HTML}{17365D}
\hypersetup{pdftitle={第076套},pdfauthor={}}
\setstretch{1.24}
\setlength{\parindent}{2em}
\setlength{\parskip}{0.35em}
\setlength{\emergencystretch}{2em}
\setlist{itemsep=0.35em,topsep=0.4em}
\captionsetup{font=small,labelfont=bf,justification=centering}
\setcounter{secnumdepth}{0}
\pagestyle{fancy}
\fancyhf{}
\fancyhead[L]{\small 第076套}
\fancyfoot[C]{\thepage}
\renewcommand{\headrulewidth}{0.4pt}
\ctexset{
  section={format=\Large\bfseries\color{ContestBlue},beforeskip=1.2em,afterskip=0.7em},
  subsection={format=\large\bfseries\color{ContestBlue},beforeskip=1.0em,afterskip=0.5em},
  subsubsection={format=\normalsize\bfseries,beforeskip=0.8em,afterskip=0.35em}
}
\usepackage{calc}
\newcounter{none}
\let\originalsubsection\subsection
\renewcommand{\subsection}{\needspace{6\baselineskip}\originalsubsection}
\makeatletter
\newcommand{\tightlist}{\setlength{\itemsep}{0pt}\setlength{\parskip}{0pt}}
\newcommand{\maxwidth}{\ifdim\Gin@nat@width>\linewidth\linewidth\else\Gin@nat@width\fi}
\newcommand{\maxheight}{\ifdim\Gin@nat@height>0.72\textheight 0.72\textheight\else\Gin@nat@height\fi}
\newcommand{\pandocbounded}[1]{\begingroup\setkeys{Gin}{width=\maxwidth,height=\maxheight,keepaspectratio}#1\endgroup}
\makeatother
\begin{document}
\begin{titlepage}
  \centering
  \vspace*{58mm}
  {\Huge\bfseries\color{ContestBlue} 第076套\par}
  \vfill
  \vspace*{28mm}
\end{titlepage}
\section{第 1 题}\label{ux7b2c-1-ux9898}

一颗恒星在形成之前，经历半径很大、密度很低的原恒星阶段。由于引力作用，原恒星逐渐收缩，释放的引力势能先转化为物质的内能，随后以辐射形式向外传递。以太阳为例，设太阳及其早期原恒星都可近似看成均匀球体。已知太阳质量
\(M=2.0\times10^{30}\,\mathrm{kg}\)、半径
\(R=7.0\times10^8\,\mathrm m\)、表面温度约为
\(T=6000\,\mathrm K\)。万有引力常量及斯特藩---玻尔兹曼常量为

\[G=6.7\times10^{-11}\,\mathrm{N\,m^2/kg^2},\qquad
\sigma=5.7\times10^{-8}\,\mathrm{W/(m^2\,K^4)}.\]

（1）太阳由半径趋于无穷大的均匀原恒星收缩为当前半径的均匀球体时，共释放多少引力势能？

（2）按照太阳当前的黑体辐射功率估算，第（1）问所得能量相当于太阳辐射多长时间？

（3）若第（1）问的能量没有被缓慢辐射，而是瞬时全部转化为太阳物质的内能，取平均比热容

\[
c=1.0\times10^3\ \mathrm{J/(kg\,K)},
\]

且忽略较低的初温，估算太阳可达到的温度。进一步假设恒星的平均密度和平均比热容都与太阳近似相同，核聚变点火温度为

\[
T_c=2.0\times10^7\ \mathrm K,
\]

估算能够达到该温度的恒星质量至少约为太阳质量的百分之几。

\begin{figure}[H]
\centering
\includegraphics[width=0.42\linewidth,height=\textheight,keepaspectratio,alt={图1-1　恒星由稀疏状态收缩并向外辐射的物理过程}]{../figures/question-01.pdf}
\caption*{图1-1　恒星由稀疏状态收缩并向外辐射的物理过程}
\end{figure}

\newpage

\section{第 2 题}\label{ux7b2c-2-ux9898}

在理想光滑水平面上放置三个相同的点电荷
\(A\)、\(B\)、\(C\)。每个点电荷的质量为 \(m\)，电荷量为
\(q\)。\(AB\)、\(AC\) 之间分别用长度为 \(l\)
的不可伸长柔软绝缘细绳连接。初始时三点构成边长为 \(l\)
的正三角形，且都由静止释放。

（1）当 \(\angle BAC=2\theta\) 时，求点电荷 \(A\) 的速度大小。

（2）求点电荷 \(A\) 的最大速度。

（3）该系统作周期运动，求运动周期 \(T\)。

\begin{figure}[H]
\centering
\includegraphics[width=0.54\linewidth,height=\textheight,keepaspectratio,alt={图2-1　三点电荷的初始构型与两根等长约束绳}]{../figures/question-02.pdf}
\caption*{图2-1　三点电荷的初始构型与两根等长约束绳}
\end{figure}

\newpage

\section{第 3 题}\label{ux7b2c-3-ux9898}

两块相同的薄金属圆盘 \(A\)、\(B\)，直径均为
\(5.0\ \mathrm{cm}\)。两圆盘用绝缘细线悬挂，彼此平行且间距为 \(d\)，其中
\(d\) 远小于圆盘直径。每块圆盘带电荷量
\(q\)，两圆盘之间的静电排斥作用明显。

现将一块不带电的大薄金属圆盘 \(C\) 小心插入 \(A\)、\(B\) 之间。圆盘
\(C\) 的直径 \(D>5.0\ \mathrm{cm}\)，不接触
\(A\)、\(B\)；三圆盘互相平行，中心位于同一直线上，且 \(A\)、\(B\) 关于
\(C\) 对称。求 \(D\) 约为多大，才能使每个圆盘所受静电力的合力均为零。

\begin{figure}[H]
\centering
\includegraphics[width=0.9\linewidth,height=\textheight,keepaspectratio,alt={图3-1　三块同轴平行圆盘的侧视关系与正视投影}]{../figures/question-03.pdf}
\caption*{图3-1　三块同轴平行圆盘的侧视关系与正视投影}
\end{figure}

\newpage

\section{第 4 题}\label{ux7b2c-4-ux9898}

\textbf{（a）} 图示电阻网络中每个电阻的阻值均为 \(R\)，求 \(A\)、\(B\)
两点之间的等效电阻。

\textbf{（b）} 图示带通滤波电路中，输入电压为

\[
V_{\rm in}=V_0\cos\omega t,
\]

其中 \(V_0\) 为定值，角频率 \(\omega\) 可调。电源经串联电阻 \(R_s\)
驱动由电容 \(C\)、电感 \(L\) 和负载电阻 \(R_L\) 组成的并联支路，输出电压
\(V_{\rm out}\) 取并联支路两端电压。

（1）当 \(\omega\) 为何值时，\(V_{\rm out}\)
的振幅最大？最大振幅是多少？

（2）当 \(R_L\) 消耗的平均功率等于其最大值的一半时，求对应的两个角频率。

（3）带通滤波器的频宽定义为第（2）问两个频率之差

\[
\Delta f=\frac{\Delta\omega}{2\pi},
\]

求 \(\Delta f\)。

\begin{figure}[H]
\centering
\includegraphics[width=0.9\linewidth,height=\textheight,keepaspectratio,alt={图4-1　对称电阻网络与并联谐振带通电路}]{../figures/question-04.pdf}
\caption*{图4-1　对称电阻网络与并联谐振带通电路}
\end{figure}

\newpage

\section{第 5 题}\label{ux7b2c-5-ux9898}

水平地面上放置由两个相同刚体组成的系统，两刚体在 \(O\)
点用光滑铰链连接。每个刚体都由半径为 \(R\) 的半圆与长为 \(L\)、宽为
\(2R\) 的矩形拼接而成，\(O\) 为半圆圆心。每个刚体的质量为
\(m\)，质心在对称轴上，且到 \(O\) 的距离为

\[
|OC|=\frac L2.
\]

刚体相对于质心的转动惯量为

\[
I_C=\frac14mL^2.
\]

上刚体顶点 \(D\) 与下刚体顶点 \(A\) 之间连接一根轻质弹性绳，弹性绳原长为
\(L\)、劲度系数为
\(k\)。绳贴着刚体表面，能在表面无摩擦滑动。开始时下刚体固定，两刚体重合，即夹角
\(\theta=\pi\)；将上刚体由静止释放。定义

\[
\alpha=\frac RL,
\qquad
\beta=\frac{mg}{kL}.
\]

（1）已知 \(\beta<2\alpha^2\)，证明上刚体能够到达最高位置 \(\theta=0\)。

（2）在第（1）问条件下，当上刚体到达最高位置时，立即解除下刚体的固定。若要求下刚体在该瞬间仍保持静止，求
\(\alpha\)、\(\beta\) 以及地面与下刚体间摩擦因数 \(\mu\) 应满足的条件。

\begin{figure}[H]
\centering
\includegraphics[width=0.597\linewidth,height=\textheight,keepaspectratio,alt={图5-1　两复合刚体、光滑铰链及贴面弹性绳的几何关系}]{../figures/question-05.pdf}
\caption*{图5-1　两复合刚体、光滑铰链及贴面弹性绳的几何关系}
\end{figure}

\newpage

\section{第 6 题}\label{ux7b2c-6-ux9898}

一个圆柱形管道内有可无摩擦移动的活塞，管内封闭 \(2.00\ \mathrm{mol}\)
空气。初态处于一个标准大气压，体积为 \(V_0\)，温度为

\[
T_0=298\ \mathrm K.
\]

气体依次经历下列过程：

\begin{itemize}
\tightlist
\item
  过程 \(A\)：在温度 \(T_0\) 下等温压缩到体积 \(V_0/4\)；
\item
  过程 \(B\)：绝热膨胀到体积 \(V_3=15.0\ \mathrm L\)；
\item
  过程 \(C\)：在过程 \(B\) 末态温度下等温膨胀到体积 \(V_0\)；
\item
  过程 \(D\)：保持体积 \(V_0\) 不变，使气体回到初温 \(T_0\)。
\end{itemize}

空气视为双原子理想气体，取

\[
C_p=\frac72R_g,
\qquad
C_v=\frac52R_g,
\qquad
1\ \mathrm{atm}=1.01\times10^5\ \mathrm{Pa}.
\]

（1）画出整个循环的 \(P\)---\(V\) 图。

（2）求过程 \(A\) 中外界对气体所做的功。

（3）求过程 \(B\) 末态气体温度。

（4）求整个循环中的最小压强。

（5）若将整个循环反向进行，求反向循环作为热机时的效率。

\begin{figure}[H]
\centering
\includegraphics[width=0.78\linewidth,height=\textheight,keepaspectratio,alt={图6-1　活塞气体的四个准静态过程}]{../figures/question-06.pdf}
\caption*{图6-1　活塞气体的四个准静态过程}
\end{figure}

\newpage

\section{第 7 题}\label{ux7b2c-7-ux9898}

在测量杨氏模量的实验中，有一根原长为 \(l\)、半径为 \(r\)、密度为
\(\rho\) 的均匀细绳。细绳与后方的接地金属刻度板平行，二者距离为
\(d\)，且刻度板足够宽、细绳与刻度板均足够长。现使电荷量 \(q\)
均匀分布在细绳外表面。实验中使用的负重平均重量为
\(T\)，按通常方法测得杨氏模量为
\(Y\)。在最终误差分析中，求静电作用给测量结果带来的修正量 \(\Delta Y\)。

可采用细绳近似，且满足

\[
r\ll d\ll l.
\]

\begin{figure}[H]
\centering
\includegraphics[width=0.622\linewidth,height=\textheight,keepaspectratio,alt={图7-1　带电细绳与接地金属刻度板的空间位置}]{../figures/question-07.pdf}
\caption*{图7-1　带电细绳与接地金属刻度板的空间位置}
\end{figure}

\newpage

\section{第 8 题}\label{ux7b2c-8-ux9898}

在量子化的一维简谐运动中，粒子的位移形式与经典简谐运动相同。设处于最低能量运动状态时

\[
x(t)=x_0\cos\omega t,
\qquad
\omega=\sqrt{\frac{k}{m}},
\]

\(x_0\) 为振幅，\(k\) 为力常量，\(m\) 为粒子质量。尖括号
\(\langle\ \rangle\) 表示一个周期内的时间平均。位移与动量不确定度定义为

\[
\Delta x=\sqrt{\left\langle(x-\langle x\rangle)^2\right\rangle},
\qquad
\Delta p=\sqrt{\left\langle(p-\langle p\rangle)^2\right\rangle}.
\]

（1）证明

\[
(\Delta x)^2=\langle x^2\rangle,
\qquad
(\Delta p)^2=\langle p^2\rangle.
\]

（2）利用测不准关系
\(\Delta x\,\Delta p\geq\hbar/2\)，证明该粒子可能的最小振幅为
\(x_{0,\min}=\sqrt{\hbar/(m\omega)}\)，并求对应的系统最小总能量。

（3）若粒子带电荷量 \(q\)，并置于沿振动方向的微弱匀强电场中，电场强度为
\(\mathcal E\)。已知粒子仍保持简谐振动特性，求最低总能量和相应振幅。

\begin{figure}[H]
\centering
\includegraphics[width=0.4\linewidth,height=\textheight,keepaspectratio,alt={图8-1　一维简谐运动的位移---时间关系}]{../figures/question-08.pdf}
\caption*{图8-1　一维简谐运动的位移---时间关系}
\end{figure}
\end{document}
